
\chapter{Controllo Ottimo Stocastico a Tempo Discreto}

In questo capitolo introduciamo la teoria del controllo ottimo stocastico a tempo discreto e illustriamo il metodo della \textbf{Programmazione Dinamica} per risolverlo.

\section{Formulazione Generale del Problema}

Molti problemi applicati (finanza, economia, gestione scorte) rientrano in questo schema generale.
Gli elementi costitutivi sono:
\begin{enumerate}
    \item \textbf{Spazio degli stati} $S \subset \R^k$: dove vive il sistema.
    \item \textbf{Spazio dei controlli} $C \subset \R^d$: le azioni a disposizione dell'agente.
    \item Un \textbf{processo stocastico} $\xi = (\xi_n)_{n \in \N}$ i.i.d., che rappresenta l'alea del sistema (il caso).
    \item La \textbf{dinamica} del sistema, data da una funzione $f$:
    \[ X_{n+1} = f(X_n, c_n, \xi_{n+1}) \]
    dove $X_n$ è lo stato attuale, $c_n$ l'azione scelta e $\xi_{n+1}$ lo shock casuale in ingresso.
    \item I \textbf{vincoli} sulle azioni ammissibili $c_n \in \mathcal{A}(x)$.
    \item Un \textbf{payoff corrente} (funzione di running cost/reward) $g(X_n, c_n)$ e un \textbf{payoff finale} $\Phi(X_N)$.
\end{enumerate}

L'obiettivo è massimizzare (o minimizzare) il valore atteso della somma dei payoff, solitamente scontati con un fattore $\beta \in (0, 1]$:
\[ \sup_{c \in \mathcal{A}(x)} \E \left[ \sum_{n=0}^{N-1} \beta^n g(X_n, c_n) + \beta^N \Phi(X_N) \right] \]

\begin{spiegazione}[Anello Aperto vs Anello Chiuso]
Il controllo $c_n$ deve essere non-anticipante, ovvero deciso basandosi solo sulle informazioni disponibili fino al tempo $n$. 
Si dice "in anello aperto" se dipende dall'intera storia passata, o "in anello chiuso" (di feedback markoviano) se $c_n = \gamma(X_n)$, cioè l'azione ottimale dipende solo dallo stato \emph{corrente}.
\end{spiegazione}

\section{Esempi Classici}

\subsection{Crescita Economica Stocastica Ottima}
Un agente (o nazione) produce una ricchezza $Q_n$ ad ogni anno. Tale ricchezza dipende da un fattore di produttività casuale (es. clima, tecnologia) $Z_n$. Ad ogni anno, l'agente sceglie quanta parte della ricchezza $c_n$ consumare. La parte non consumata viene reinvestita e produce la ricchezza dell'anno successivo. L'obiettivo è massimizzare l'utilità attesa dei consumi nel tempo: $\E \sum \beta^n U(c_n)$.

\subsection{Allocazione di Portafoglio}
Hai un conto in banca a tasso privo di rischio $r$ e un asset rischioso (es. azioni) i cui ritorni aleatori sono $\mu_n$. Il tuo stato è il valore totale del portafoglio $W_n$. Il controllo $\pi_n$ è la frazione di ricchezza che investi nel titolo rischioso. L'obiettivo è massimizzare l'utilità della ricchezza finale: $\E [ U(W_N) ]$.

\subsection{Gestione Scorte (Inventory)}
Un magazzino deve rifornirsi ordinando $c_n$ unità di merce ad ogni intervallo. Esiste una domanda casuale $D_n$. C'è un costo di acquisto per $c_n$ e un costo di penalità se le scorte finali $Q_n$ si discostano troppo dalla domanda. L'obiettivo è minimizzare i costi attesi totali.

\section{Programmazione Dinamica (DP)}

L'idea geniale della Programmazione Dinamica (di Richard Bellman) è di "rompere" il problema su tutto l'orizzonte in una sequenza di problemi a singolo passo, procedendo a ritroso nel tempo.

\subsection{La Funzione Valore e l'Equazione di Bellman}

Definiamo la \textbf{Funzione Valore} $V_k(x)$ come il massimo payoff atteso che si può ottenere \emph{se al tempo $k$ ci troviamo nello stato $x$}.

\begin{theorem}[Principio di Ottimalità di Bellman]
La funzione valore soddisfa la seguente relazione ricorsiva (Equazione di Bellman):
\[ V_k(x) = \sup_{\gamma \in \Gamma(x)} \left\{ g(x, \gamma) + \beta \E \left[ V_{k+1}(f(x, \gamma, \xi_{k+1})) \right] \right\} \]
\end{theorem}

\begin{hint}[Come Leggere Bellman]
Il significato intuitivo è: il valore ottimale al tempo $k$ si ottiene scegliendo la mossa $\gamma$ che massimizza la somma tra il guadagno immediato oggi ($g(x, \gamma)$) e il valore atteso che avrò da domani in poi, considerando che mi ritroverò nel nuovo stato generato da $\gamma$ e dal caso.
\end{hint}

\subsection{Algoritmo in Orizzonte Finito}
Se il tempo $N$ è finito, il problema si risolve partendo dalla fine.
\begin{enumerate}
    \item Si pone $V_N(x) = \Phi(x)$ (al momento terminale, incasso semplicemente il payoff finale).
    \item Per $k = N-1, N-2, \dots, 0$, si calcola $V_k(x)$ applicando l'Equazione di Bellman, supponendo noto $V_{k+1}$.
    \item Al contempo, la mossa ottimale $\gamma^*_k(x)$ è quella che realizza il sup nell'equazione di Bellman.
\end{enumerate}

\subsection{Orizzonte Infinito e Teorema di Verifica}
Se $N = \infty$, il tempo scompare ($V_k(x) = V(x)$ per ogni $k$, dato che la dinamica è stazionaria). L'Equazione di Bellman diventa una equazione a punto fisso per una singola funzione $V(x)$:
\[ V(x) = \sup_{\gamma \in \Gamma(x)} \left\{ g(x, \gamma) + \beta \E \left[ V(f(x, \gamma, \xi_1)) \right] \right\} \]

Poiché procedere a ritroso dall'infinito è impossibile, in questo caso si usa un approccio \emph{Guess \& Verify}:
\begin{enumerate}
    \item "Indovino" una funzione $v_0(x)$ che risolve l'equazione di Bellman.
    \item Trovo la strategia di feedback ottima $\gamma^*(x)$.
    \item Verifico una condizione di \emph{trasversalità} (o crescita controllata all'infinito): $\lim_{N \to \infty} \E[v_0(X_N)] \le 0$ (o equazione simile, a seconda dei segni del payoff).
    \item Il \textbf{Teorema di Verifica} assicura allora che $v_0(x)$ è \emph{la} vera funzione valore e $\gamma^*$ è il vero controllo ottimo.
\end{enumerate}

\begin{esempio}[Esempio di Orizzonte Infinito]
Nel problema di crescita (Sect 8.4.1), se supponiamo che la funzione valore abbia la forma $v_0(x) = v_0(q, z) = \alpha(z) \ln q$, sostituendola nell'equazione di Bellman otteniamo i coefficienti ottimi, trovando così analiticamente sia il valore atteso massimo sia la porzione di ricchezza esatta da consumare ogni anno in base allo stato.
\end{esempio}
